\section{ODEs}

Here, we have some basic ways to solve ODEs. We first go 
through some basic linear ODEs, which I wish to be familiar 
with and make it just an easy computational exercise.

\begin{example}[First-order linear ODE; Integrating Factors]
Consider
\[
y' + p(x)y = q(x).
\]
Define the integrating factor
\[
I(x)=e^{\int p(x) dx}.
\]
This gives
\[
y(x)
=
e^{-\int p(x) dx}
\left(
\int e^{\int p(x) dx}q(x)dx+C
\right) = \frac{1}{I(x)} \left( I(x) q(x) dx + C \right).
\]
Note that this comes from multiplying the ODE by $I(x)$ and
integrating the differentiated chain rule.
\end{example}


\begin{exercise}
Let $y(t)$ be the solution to the differential equation
\[
y' + 2ty = e^{-t^2}\sin t
\]
such that $y(0)=0$. What is $y(\pi)$?
\end{exercise}
\begin{solution}
This is first order linear. The integrating factor is $e^{t^2}$.
Thus,
\[
y(t) = e^{-t^2} \left(\int e^{t^2} \cdot e^{-t^2}\sin t dt +C \right)
= e^{-t^2} (-\cos(t) + C).
\]
$y(0) = 0$, so $C = 1$. Thus, $y(\pi) = 2e^{-\pi^2}$.
\end{solution}

\begin{example}[Second-order linear ODE with Constant Coefficients]
Consider the homogeneous equation
\[
a_2y''+a_1y'+a_0y=0.
\]
We consider $y = e^{\lambda x}$. 
Substitution gives the characteristic equation
\[
a_2\lambda^2+a_1\lambda+a_0=0.
\]
If the characteristic roots are distinct $\lambda_1 \neq \lambda_2$, then
\[
y(x)=c_1e^{\lambda_1x}+c_2e^{\lambda_2x}.
\]
If the roots coincide, say to $\lambda$, then
\[
y(x) = c_1 e^{\lambda x} + c_2 x e^{\lambda x}.
\]
If the roots are complex, so $\lambda=\alpha\pm i\beta$, 
then the real-valued general solution is
\[
y(x) = e^{\alpha x}
\left(
c_1\cos(\beta x)+c_2\sin(\beta x)
\right).
\]
\end{example}

\begin{example}[Euler-Cauchy Equation]
A $k$th-order Euler-Cauchy equation has the form
\[
\beta_k x^k \frac{d^k y}{dx^k}
+
\beta_{k-1}x^{k-1}\frac{d^{k-1}y}{dx^{k-1}}
+\cdots+
\beta_1x\frac{dy}{dx}
+
\beta_0y
=
f(x),
\]
where $\beta_0,\ldots,\beta_k$ are constants. 
The key feature is that the coefficient of the $j$th derivative 
is proportional to $x^j$.
For this, use $x = e^z$. This makes
\[
x^j\frac{d^j y}{dx^j}
\]
to be expressed as a linear combination of
\[
\frac{d^j y}{dz^j},
\frac{d^{j-1}y}{dz^{j-1}},
\ldots,
\frac{dy}{dz},
\]
with constant coefficients, so the substitution $x=e^z$ transforms 
the Euler--Cauchy equation into a linear ODE with constant coefficients 
in the variable $z$. After solving the resulting constant-coefficient ODE, substitute
$z=\ln x$ to write the solution in terms of $x$.
\end{example}

We see how this plays out below.

\begin{exercise}
Find a solution to the ODE
\[
x^2y''+ 3xy'+ y= x^3.
\]
\end{exercise}
\begin{solution}
We use $x = e^z$. We calculate the differentials first.
\[
\Dfrac{y}{x} = \Dfrac{y}{z} \Dfrac{z}{x}
= \frac{1}{x} \Dfrac{y}{z}
\]
Next, we differentiate this to find $\frac{d^2y}{dx^2}$.
\[
\frac{d^2y}{dx^2} = \Dfrac{}{x} \left(\frac{1}{x} \Dfrac{y}{z}\right) 
= -\frac{1}{x^2} \Dfrac{y}{z} + \frac{1}{x^2} \frac{d^2y}{dz^2}.
\]
Here, we used the chain rule:
\[
\Dfrac{}{x} \left(\Dfrac{y}{z}\right) = \Dfrac{}{z} \left(\Dfrac{y}{z}\right)\Dfrac{z}{x}
= \frac{1}{x^2} \frac{d^2y}{dz^2}
\]
We now substitute back in, which gives the following ODE:
\[
\frac{d^2y}{dz^2} + 2 \frac{dy}{dz} + y = e^{3z}.
\]
This is in our realm now. We attain the following:
\[
c_1 e^{-z} + c_2 z e^{-z} + \frac{1}{16} e^{3z}
\]
We just substitute back $z = \ln (x)$ to get the desired 
form for $y(x)$.
\end{solution}